\documentclass[12pt,a4paper]{article}
\usepackage[utf8]{inputenc}
\usepackage[portuguese]{babel}
\usepackage[T1]{fontenc}
\usepackage{geometry}
\usepackage{enumitem}
\usepackage{titlesec}

\geometry{margin=2.5cm}
\titleformat{\section}{\normalfont\Large\bfseries}{\thesection}{1em}{}[\titlerule]

\title{Especificação de Caso de Estudo: Projeto Rei Julian \\ \large Rastreamento Geográfico para Delivery Local}
\author{Hiago Dona Muniz \\ \small Disciplina de TEC IV - Ciência da Computação (UFPel)}
\date{22 de Abril de 2026}

\begin{document}

\maketitle

\section{Introdução}
O projeto \textbf{Rei Julian} é um sistema de monitoramento em tempo real focado na logística de entregas. O objetivo é conectar o restaurante, o motoboy e o cliente final em uma única plataforma onde a localização do pedido é compartilhada de forma automática e transparente.

\section{Importância e Impacto Prático}
A implementação desta ferramenta foca em resolver gargalos operacionais comuns no delivery:
\begin{itemize}
    \item \textbf{Eficiência no Atendimento:} Ao disponibilizar o rastreio para o cliente, o restaurante reduz drasticamente o número de chamadas telefônicas para perguntar o status da entrega.
    \item \textbf{Gestão da Frota:} O administrador tem uma visão clara da posição de cada entregador, permitindo identificar atrasos, motoboys parados ou necessidade de redistribuição de pedidos.
    \item \textbf{Experiência do Usuário:} A visibilidade do trajeto gera confiança no serviço e reduz a ansiedade do cliente durante o tempo de espera.
    \item \textbf{Independência Tecnológica:} Permite que estabelecimentos locais gerenciem sua própria logística sem depender exclusivamente de plataformas de terceiros e suas taxas.
\end{itemize}

\section{Atores do Sistema}
\begin{description}
    \item[Restaurante:] Interface administrativa para monitorar todos os entregadores ativos, gerenciar o status da frota e consultar o histórico de rotas para auditoria.
    \item[Motoboy:] Aplicativo/Interface que utiliza o GPS do smartphone para transmitir as coordenadas para o servidor e atualizar o status do pedido (Coletado, Em Rota, Entregue).
    \item[Cliente:] Link temporário de visualização onde o cliente acompanha apenas o entregador vinculado ao seu pedido específico.
\end{description}

\section{Viabilidade Técnica}
O projeto é tecnicamente viável e de baixo custo devido às escolhas arquiteturais:
\begin{itemize}
    \item \textbf{Linguagem Go:} Alta performance para lidar com múltiplas conexões simultâneas consumindo poucos recursos de servidor.
    \item \textbf{Protocolo SSE (Server-Sent Events):} Método leve para enviar atualizações de posição do servidor para o mapa, economizando bateria nos dispositivos móveis.
    \item \textbf{Persistência com SQLite:} Banco de dados eficiente que dispensa infraestruturas complexas para o armazenamento do histórico de movimentação.
    \item \textbf{Ferramentas Open Source:} Uso de Leaflet e OpenStreetMap, eliminando custos com licenciamento de mapas.
\end{itemize}

\section{Fluxo de Funcionamento}
1. O restaurante atribui uma entrega a um motoboy. 2. O sistema gera um ID de rastreio único. 3. O motoboy inicia a rota e o sistema passa a transmitir suas coordenadas. 4. O cliente visualiza o deslocamento em tempo real. 5. Ao confirmar a entrega no destino, o rastreio é encerrado e os dados são salvos para consultas futuras.

\section{Conclusão}
O Rei Julian demonstra ser uma solução prática e de fácil implementação. Sua importância está na melhoria direta da comunicação e no controle operacional do delivery, enquanto sua viabilidade é garantida pelo uso de tecnologias modernas, leves e gratuitas.

\end{document}
